\begin{definition}
	\textit{Функцией Жуковского} называется функция $f_*: \CM \to \CM$, имеющая следующий вид:
	\[f_*(z) := \System{
		& \frac12\left(z + \frac1z\right), & \text{если } &z \ne 0\text{ и } z \ne \infty
		\\
		& \infty, & \text{если } &z = 0\text{ или } z = \infty
	}\]
\end{definition}

Выделим несколько свойств функции Жуковского:

\begin{enumerate}
    \item Непрерывна в $\CM$
    \item Проверим, что функция $f_*$ конформна в каждой точке из $\CM \bs \{\pm 1\}$:
	\begin{itemize}
		\item Точки $0$ и $\infty$ являются полюсами первого порядка функции $f_*$
		\item Для любого $z \in \Cm \bs \{0\}$ функция $f_*$ регулярна в точке $z$, причем $f_*'(z) = \frac12\left(1 - \frac1{z^2}\right)$, поэтому $f_*'(z) = 0 \lra z = \pm1$
	\end{itemize}

	Кроме того, для любых $z_1, z_2 \in \CM$ условие $f_*(z_1) = f_*(z_2)$ равносильно тому, что либо $\{z_1, z_2\} = \{0, \infty\}$, либо $z_1z_2 = 1$. Это и означает требуемое.
    \item Функция $f_*$ конформна на области $D \subset \CM \Leftrightarrow D$ не содержит одновременно точки $0$ и $\infty$ и не содержит пары точек $z_1, z_2 \in \Cm$ такие, что $z_1z_2 = 1$.
\end{enumerate}

\begin{unnecessary}
Из доказательства выше следует, что функция $f_*$ конформна на любой своей области однолистности. Такими областями, в частности, являются следующие области:
\begin{itemize}
    \item $D_1 := \{z \in \Cm : |z| < 1\}$ и $D_2 := \{w \in \Cm : |w| > 1\}$, причем область $D_2$ является образом области $D_1$ при отображении $w(z) = \frac 1z$
    \item $D_3 := \{z \in \Cm : \im{z} > 0\}$ и $D_4 := \{w \in \Cm : \im{w} < 0\}$, причем область $D_4$ является образом области $D_3$ при отображении $w(z) = \frac 1z$
\end{itemize}

\begin{example}
$f_*$ конформно в области
\begin{itemize}
    \item $D_1 = \{z \in \C : |z| > 1\} \cup \{\infty \}$
    \item $D_2 = \{z \in \C : Im z > 0\}$
\end{itemize}
\end{example}
\end{unnecessary}

\begin{proposition}
	Пусть $\Omega_1 \subset \CM$ "--- область, и $f_*$ конформна в этой области, $\Omega_2$ "--- область, являющаяся образом области $\Omega_1$ при отображении $w(z) = \frac 1z$. Тогда $f_*$ конформна в $\Omega_2$
\end{proposition}

\begin{proof}
    Т.к. $\Omega_1$ получается из $\Omega_2$ конформным отображением $w = \frac{1}{z}$, то отображение $g = f_*(w)$ конформно в $\Omega_2$ как композиция конформных отображений. Но $f_*(\frac{1}{z}) = f_*(z)$. Значит, $f_*(z)$ конформна в $\Omega_2$
\end{proof}

\begin{proposition}
Отображение окружности и лучей:

\begin{itemize}
    \item $f_*(D_1) = f_*(D_2) = \CM \bs [-1, 1]$
    \item $f_*(D_3) = f_*(D_4) = \CM \bs \big((-\infty, -1] \cup [1, +\infty)\big)$
    \item Окружность с центром в точке $0$ и радиусом $\rho \ne 1$ под действием функции Жуковского переходит в эллипс с фокусами $\pm 1$, задаваемый следующим уравнением:
    \[\frac{u^2}{\left(\frac12\left(\rho + \frac 1\rho\right)\right)^2} + \frac{v^2}{\left(\frac12\left(\rho - \frac 1\rho\right)\right)^2} = 1\]
    \item Открытый луч с началом в точке $0$, состоящий из точек с аргументом $\phi \in \left(0, \frac\pi 2\right)$, под действием функции Жуковского переходит в ветвь гиперболы, задаваемой следующим уравнением:
    \[\frac{u^2}{\cos^2\phi} - \frac{v^2}{\sin^2\phi} = 1\]
\end{itemize}
\end{proposition}
